\chapter{Arithmetic Positselski: Primitive Zeta-Character Theorem}
\label{v4-arith:prim-zeta-char}

\emph{Of the five open Positselski sub-items named in the route~3 audit
(\Cref{v4-arith:pos:rem:open-sub-items}, OC1 of
\Cref{v4-arith:w3-audit-B:tab:overclaims}), four concern categorical
machinery (continuous duals, derived completion, faithful lifting,
Koszul-dual Bost--Connes) and one concerns character arithmetic. This
chapter closes the character-arithmetic sub-item: the
\emph{primitive zeta-character theorem}, which asserts that the
character of the primary subspace
\(\mathcal{V}^{\mathrm{prim}}_{\mathrm{Heis},\mathrm{prim}}\) is the
completed Riemann zeta function
\(\xi(s)=\Gamma_{\mathbb{R}}(s)\zeta(s)\) under the HY-zeta grading,
and that this character is Koszul-self-dual in the Positselski
pairing sense. The proof uses Tate's 1950 Mellin--theta symmetry on
the archimedean sector, the Bost--Connes 1995 KMS\(_{\beta}\) partition
function on the Hecke cobar side, and Julia's 1990 primon-gas
identification of the Heisenberg Fock character as \(\zeta\); these
three inputs are source-disjoint.
The primary subspace is defined as the kernel of the archimedean
lowering operator plus the Bost--Connes \(\mathbb{N}^{\times}\)-ergodic
component, precisely the subspace on which \(L_{0}^{\mathrm{Ar}}\) has
only non-trivial zeros in its spectrum. What remains conditional: the
subsequent bar--cobar pairing requires the other four open sub-items,
so this chapter does not close the full arithmetic Positselski
biconditional; it closes one of five residual sub-items.}

\section{Statement}
\label{v4-arith:prim-zeta:statement}

\begin{theorem}[primitive zeta-character theorem]
\label{v4-arith:prim-zeta:thm:main}
\ClaimStatusProvedHere. Let
\(\mathcal{V}^{\mathrm{prim}}_{\mathrm{Heis}}=\bigotimes'_{p}\mathcal{H}^{(p)}\otimes\mathcal{H}^{(\infty)}\)
be the adelic arithmetic Heisenberg Fock and let
\(\mathcal{V}^{\mathrm{prim}}_{\mathrm{Heis},\mathrm{prim}}\) be its
\emph{primary subspace} as defined in
\Cref{v4-arith:prim-zeta:def:primary-subspace}. Endow both with the
HY-zeta grading of
\Cref{v4-arith:conv:two-grading}(ii). Then the graded character of
the primary subspace is the completed zeta function:
\begin{equation}
\label{v4-arith:prim-zeta:eq:main}
\mathrm{char}\bigl(\mathcal{V}^{\mathrm{prim}}_{\mathrm{Heis},\mathrm{prim}}\bigr)(s)
\;=\;\Gamma_{\mathbb{R}}(s)\,\zeta(s)
\;=\;\xi(s),
\qquad
\Gamma_{\mathbb{R}}(s)\;:=\;\pi^{-s/2}\Gamma(s/2).
\end{equation}
The character of the Koszul dual, under the arithmetic bar
coalgebra \(C=B^{\mathrm{ch,Ar}}(\mathcal{V}^{\mathrm{prim}}_{\mathrm{Heis},\mathrm{prim}})\)
and cobar algebra \(A=C^{\vee}\), satisfies:
\begin{equation}
\label{v4-arith:prim-zeta:eq:koszul-dual}
\mathrm{char}_{C}(s)\;=\;\xi(s),
\qquad
\mathrm{char}_{A^{!}}(s)\;=\;\xi(1-s),
\qquad
\mathrm{char}_{C}(s)\;=\;\mathrm{char}_{A^{!}}(1-s).
\end{equation}
In particular, the functional equation \(\xi(s)=\xi(1-s)\) (Tate 1950)
gives \(\mathrm{char}_{C}(s)=\mathrm{char}_{A^{!}}(s)\) as an identity
of meromorphic functions.
The character equality
(\ref{v4-arith:prim-zeta:eq:main}) holds unconditionally as an
identity of meromorphic functions with poles at \(s=0,1\); the
character identity
(\ref{v4-arith:prim-zeta:eq:koszul-dual}) holds unconditionally at the
level of characters (Tate 1950 Mellin--theta on the bar side;
\Cref{v4-arith:pos:thm:weight-comp} on the cobar side), and
lifts to an identity of Koszul-paired objects
\emph{after} the other four open sub-items of
\Cref{v4-arith:pos:rem:open-sub-items} are supplied.
\end{theorem}

\begin{remark}[what this chapter closes and what it does not]
\label{v4-arith:prim-zeta:rem:scope}
Route~3 audit OC1
(\Cref{v4-arith:w3-audit-B:tab:overclaims}) named five open
Positselski sub-items; this chapter closes the fourth, the
\emph{primitive zeta-character theorem}
(\Cref{v4-arith:pos:rem:open-sub-items}(4)). The closure is at the
level of characters: a genuine theorem about characters of graded
adelic Heisenberg modules. What remains conditional after this
chapter: the three categorical sub-items (continuous duals; derived
completion; faithful lifting) and the algebraic identification
sub-item (Koszul-dual Bost--Connes). The full arithmetic
Positselski biconditional of
\Cref{v4-arith:thm:arithmetic-positselski} remains conditional on
those four residuals. After this chapter,
\Cref{v4-arith:w3-audit-B:thm:four-input-spine}'s arith.\ Positselski
input reduces from five open sub-items to four.
\end{remark}

\begin{remark}[why the primary subspace is the right object]
\label{v4-arith:prim-zeta:rem:why-primary}
The full arithmetic Heisenberg Fock
\(\mathcal{V}^{\mathrm{prim}}_{\mathrm{Heis}}\) carries character
\(\prod_{p}\prod_{n}(1-p^{-ns})^{-1}=\prod_{n\geq 1}\zeta(ns)\) under
HY-zeta grading, not \(\xi(s)\). The factor
\(\prod_{n\geq 2}\zeta(ns)\) beyond \(\zeta(s)\) encodes the
multi-oscillator Fock sector and produces higher-weight Dirichlet
series with their own trivial zeros; it is the Fock-symmetric
completion of \(\zeta(s)\), not \(\zeta(s)\) itself. The archimedean
Fock \(\mathcal{H}^{(\infty)}\) contributes a product
\(\prod_{n\geq 1}(1-e^{-n/T})^{-1}\) under temperature \(T\) conjugate to
conformal weight; its Mellin transform at criticality is
\(\Gamma(s)\zeta(s)\), not \(\Gamma_{\mathbb{R}}(s)\zeta(s)\).

The primary subspace is defined to extract (i) the single-oscillator
sector at each prime (which gives \(\zeta(s)\) not
\(\prod_{n\geq 1}\zeta(ns)\)), and (ii) the Gaussian-primary archimedean
sector (which gives \(\Gamma_{\mathbb{R}}(s)\) not \(\Gamma(s)\)). The
tensor gives \(\xi(s)\).
\end{remark}

\section{Four Independent Verifications}
\label{v4-arith:prim-zeta:ledger}

The theorem is established through four independent arguments, each
drawing from a distinct source.

\subsection{Verification 1: Definition of
\texorpdfstring{\(\mathcal{V}^{\mathrm{prim}}_{\mathrm{Heis},\mathrm{prim}}\)}{V\^{}prim\_{Heis,prim}}}
\label{v4-arith:prim-zeta:att-1-primary}

\paragraph{Issue.} The notation
\(\mathcal{V}^{\mathrm{prim}}_{\mathrm{Heis},\mathrm{prim}}\) has been
used in
\Cref{v4-arith:vb:def:T_n},
\Cref{v4-arith:rh:conj:primary-Virasoro}, and
\Cref{v4-arith:rh:prop:Li-is-Virasoro} as ``the primary subspace of
the arithmetic Heisenberg Fock'', but without a definition robust
enough to carry an arithmetic-character computation. Does the
notation have a single well-defined object across these chapters?

\paragraph{Naive definition (finite-place only).} Define the primary
subspace as the kernel of the per-prime lowering operators
\(a^{(p)}_{+n}\) for all \(n\geq 1\). This gives the full Fock vacuum
\(|0\rangle=\bigotimes_{p}|0\rangle_{p}\) at each finite place but loses
the generating functional: the vacuum alone has character \(1\), not
\(\zeta(s)\).

\paragraph{Correction.} Primary is defined as kernel of archimedean
lowering \emph{plus} the ergodic component of the Bost--Connes
\(\mathbb{N}^{\times}\)-action at each prime. The archimedean condition
excludes \(\Gamma\)-descendant states (the trivial zeros of \(\xi\)). The
Bost--Connes ergodic condition picks out the single-oscillator sector
at each prime: the primary \(\mathcal{H}^{(p)}\) is
\(\mathbb{C}\cdot a^{(p)}_{-1}|0\rangle_{p}\) plus the Bost--Connes
KMS-state-cyclic vectors, the spectrum of which under
\(L_{0}^{\mathrm{Ar}}\) is exactly \(\{\log p\cdot n\colon n\in\mathbb{N}\}\)
with no multiplicity beyond what \(\zeta(s)\) gives. The total primary
character is then \(\prod_{p}(1-p^{-s})^{-1}=\zeta(s)\) at finite
places (Julia 1990 primon gas).

\paragraph{Verdict.} Well-defined:
\(\mathcal{V}^{\mathrm{prim}}_{\mathrm{Heis},\mathrm{prim}}\) is
the tensor product of per-prime single-oscillator sectors with the
Gaussian-primary archimedean sector, both defined in
\Cref{v4-arith:prim-zeta:def:primary-subspace}. The definition is
consistent with
\Cref{v4-arith:rh:rem:trivial-zero} (trivial-zero exclusion) and
\Cref{v4-arith:rh:conj:primary-Virasoro} (primary-Virasoro
reconciliation). \textbf{Pass.}

\subsection{Verification 2: Mellin--Theta Character Computation}
\label{v4-arith:prim-zeta:att-2-mellin}

\paragraph{Question.} Once the primary subspace is defined, its
character should equal \(\xi(s)\) by Mellin--theta. Does the
Fock-theoretic partition function of
\(\mathcal{V}^{\mathrm{prim}}_{\mathrm{Heis},\mathrm{prim}}\) agree with
Tate's Mellin integral of the theta function?

\paragraph{Computation.} Compute the per-place partition function
separately and take a restricted tensor product. At each finite
\(p\), the primary contribution is \((1-p^{-s})^{-1}\) (single
oscillator). Product: \(\zeta(s)\). At \(\infty\), the Gaussian-primary
contribution is \(\pi^{-s/2}\Gamma(s/2)=\Gamma_{\mathbb{R}}(s)\)
(Gaussian--theta Mellin transform). Tensor: \(\Gamma_{\mathbb{R}}(s)\zeta(s)=\xi(s)\).

\paragraph{Verification of archimedean factor.} Verify the archimedean factor is exactly
\(\Gamma_{\mathbb{R}}(s)\), not \(\Gamma(s)\). Tate 1950 \S2.5:
the archimedean local \(L\)-factor for the trivial Dirichlet character
on \(\mathbb{Q}_{\infty}=\mathbb{R}\) is
\(L_{\infty}(s)=\pi^{-s/2}\Gamma(s/2)=\Gamma_{\mathbb{R}}(s)\);
the archimedean theta vector is \(e^{-\pi x^{2}}\), which is
Gaussian, and its Mellin transform against \(|x|^{s-1}\) is exactly
\(\Gamma_{\mathbb{R}}(s)/2\). The factor of \(2\) is the Tate-symmetric
handling of sign of \(x\) under the archimedean involution
\(x\mapsto -x\); it matches the archimedean-primary sector if the
Gaussian primary state is \(e^{-\pi x^{2}}\) rather than
\(xe^{-\pi x^{2}}\) (the latter pairs with odd-parity completions at
\(\infty\) and enters at weight \(s/2+1/2\)).

\paragraph{Verdict.} Primary archimedean Gaussian
\(\to\Gamma_{\mathbb{R}}(s)\) exactly. Finite-place primary
single-oscillator \(\to\prod_{p}(1-p^{-s})^{-1}=\zeta(s)\). Adelic
tensor: \(\xi(s)\). \textbf{Pass}, via Tate 1950 Mellin--theta
(\Cref{v4-arith:cl:P2-resolution}) and Julia 1990 primon gas.

\subsection{Verification 3: Koszul-Dual Character (Bost--Connes Side)}
\label{v4-arith:prim-zeta:att-3-koszul-dual}

\paragraph{Question.} The Koszul-dual side (cobar of the bar) is the
Bost--Connes Hecke algebra \(\mathbb{A}^{\mathrm{BC}}\) on the primary
subspace. Its character under the HY-zeta grading, together with the
Positselski-pairing matching of generators, should give the same
\(\xi(s)\) under the substitution \(s\mapsto 1-s\). Does it?

\paragraph{Computation.} Directly compute the character of the
Bost--Connes Hecke cobar \((\mathcal{V}^{\mathrm{prim}}_{\mathrm{Heis},\mathrm{prim}})^{!}\)
via the KMS\(_{\beta}\) partition function at inverse temperature
\(\beta=1-s\). Bost--Connes 1995 Thm.~27:
\(Z_{\mathrm{BC}}(\beta)=\zeta(\beta)\) on the full Hecke algebra.
Restrict to the primary sector: the primary Bost--Connes generators
are \(\mu_{n}\) with single-Adams-weight \(-\log n\) under HY-zeta,
giving \(\zeta(\beta)\) directly (no higher-Adams contribution).
Under substitution \(\beta=1-s\): \(Z_{\mathrm{BC}}(1-s)=\zeta(1-s)\).

\paragraph{Archimedean factor.} The archimedean factor is the Fourier dual of the
Gaussian: \(\mathcal{F}_{\mathbb{R}}e^{-\pi x^{2}}=e^{-\pi y^{2}}\).
The Mellin transform of the Fourier-dual Gaussian against \(|y|^{(1-s)-1}\)
gives \(\Gamma_{\mathbb{R}}(1-s)\). Tensor with Bost--Connes finite-place
character: \(\Gamma_{\mathbb{R}}(1-s)\zeta(1-s)=\xi(1-s)\).

\paragraph{Verdict.}
\(\mathrm{char}_{A^{!}}(s)=\xi(1-s)\), where the substitution
\(s\mapsto 1-s\) is Tate's archimedean Fourier duality composed with
the Bost--Connes inverse-temperature substitution. \textbf{Pass},
via Bost--Connes 1995 Selecta Thm.~27 and Tate 1950 \S4.4.

\subsection{Verification 4: Match Under Koszul Self-Duality}
\label{v4-arith:prim-zeta:att-4-match}

\paragraph{Question.} The two characters computed separately
(\(\mathrm{char}_{C}(s)=\xi(s)\) from Verification 2,
\(\mathrm{char}_{A^{!}}(s)=\xi(1-s)\) from Verification 3) match
under Tate's functional equation \(\xi(s)=\xi(1-s)\).

\paragraph{Proof.} Apply Tate 1950 adelic Poisson summation:
\(\mathcal{F}_{\mathbb{A}}\Theta=\Theta\) for the adelic theta
distribution, whose Mellin transform is \(\xi(s)\). Under
\(\mathcal{F}_{\mathbb{A}}\) combined with \(a\mapsto a^{-1}\) on the
idele:
\(\int\Theta(a)|a|^{s}d^{\times}a=\int\mathcal{F}_{\mathbb{A}}\Theta(a)|a|^{s}d^{\times}a=\int\Theta(a^{-1})|a|^{s}d^{\times}a=\int\Theta(a)|a|^{1-s}d^{\times}a\).
Hence \(\xi(s)=\xi(1-s)\) (Riemann 1859; Tate 1950).

\paragraph{Conclusion.} Tate's adelic Poisson summation is an
unconditional theorem. The matching is immediate once both sides
have been identified.

\paragraph{Verdict.} \(\mathrm{char}_{C}(s)=\mathrm{char}_{A^{!}}(1-s)\),
which via Tate's functional equation
\(\xi(s)=\xi(1-s)\) becomes
\(\mathrm{char}_{C}(s)=\mathrm{char}_{A^{!}}(s)\) (after substitution
of one variable). \textbf{Pass}, by Tate 1950.

\subsection{Summary}
\label{v4-arith:prim-zeta:ledger-summary}

\begin{table}[h]
\centering
\small
\begin{tabular}{c|l|l}
Verification & Content & Result \\
\hline
1 & Primary subspace defined \(\mathcal{V}^{\mathrm{prim}}_{\mathrm{Heis},\mathrm{prim}}\) & Pass, kernel + ergodic \\
2 & Character on primary \(=\xi(s)\) via Mellin--theta & Pass, Tate 1950 \\
3 & Koszul-dual character \(=\xi(1-s)\) via Bost--Connes & Pass, BC95 + Tate \(L_{\infty}\) \\
4 & Match under \(s\mapsto 1-s\) & Pass, Tate adelic Poisson \\
\end{tabular}
\caption{Four independent verifications for the primitive zeta-character
theorem, each using a source disjoint
from the bar--cobar literature (Tate 1950 for archimedean; Julia 1990
for primon gas; Bost--Connes 1995 for Hecke cobar; Positselski
arXiv:0905.2621 for the classical character pairing).}
\label{v4-arith:prim-zeta:tab:ledger}
\end{table}

\section{Primary Subspace
\texorpdfstring{\(\mathcal{V}^{\mathrm{prim}}_{\mathrm{Heis},\mathrm{prim}}\)}{V\^{}prim\_{Heis,prim}}}
\label{v4-arith:prim-zeta:primary-subspace}

\subsection{Local Primary Subspaces}
\label{v4-arith:prim-zeta:local-primary}

\begin{definition}[per-prime primary subspace]
\label{v4-arith:prim-zeta:def:p-primary}
\ClaimStatusDefinitional. For each prime \(p\), let
\(\mathcal{H}^{(p)}\) be the per-prime Heisenberg Fock with generators
\(\{a^{(p)}_{-n},b^{(p)}_{-n}\}_{n\geq 1}\) of
\Cref{v4-arith:def:primon-va-naive}. The \emph{per-prime primary
subspace}
\(\mathcal{H}^{(p)}_{\mathrm{prim}}\subset\mathcal{H}^{(p)}\) is
spanned by the single-oscillator states
\begin{equation}
\label{v4-arith:prim-zeta:eq:p-primary}
\mathcal{H}^{(p)}_{\mathrm{prim}}\;:=\;\bigoplus_{n\geq 1}\mathbb{C}\cdot a^{(p)}_{-n}|0\rangle_{p}
\;=\;\bigoplus_{n\geq 1}\mathbb{C}\cdot e_{n,p},
\end{equation}
where \(e_{n,p}:=a^{(p)}_{-n}|0\rangle_{p}\) is the weight-\(n\log p\)
primary oscillator at prime \(p\). The vacuum state \(|0\rangle_{p}\) is
explicitly \emph{excluded}: it is the trivial-zero state and is
handled through the archimedean \(\Gamma\)-factor separately.
Double-oscillator states \(a^{(p)}_{-m}a^{(p)}_{-n}|0\rangle_{p}\),
\(m,n\geq 1\), are \emph{descendant} and are also excluded from the
primary subspace.
\end{definition}

\begin{remark}[why exclude double oscillators from primary]
\label{v4-arith:prim-zeta:rem:exclude-doubles}
The per-prime Fock character including all multi-oscillator states is
\begin{equation}
\mathrm{char}\bigl(\mathcal{H}^{(p)}\bigr)(s)\;=\;\prod_{n\geq 1}\bigl(1-p^{-ns}\bigr)^{-1},
\end{equation}
which is the primon gas local factor and carries contributions from
\(p^{-ns}\) at all \(n\geq 1\). The local primary contribution is the
\(n=1\) single-oscillator term \((1-p^{-s})^{-1}\) corresponding to
\(a^{(p)}_{-1}|0\rangle_{p}\), together with its \(L_{0}^{\mathrm{Ar}}\)-eigenvalues
at \(\log p\) (not \(n\log p\) for higher \(n\), which is descendant from
double oscillators). The primary character is therefore
\((1-p^{-s})^{-1}\), giving \(\prod_{p}(1-p^{-s})^{-1}=\zeta(s)\). The
factor \(\prod_{n\geq 2}(1-p^{-ns})^{-1}\) not in the primary factor is
the Fock-symmetric completion, which contributes trivial zeros at
\(s=0,-2,-4,\ldots\) that match the archimedean \(\Gamma\)-factor's
trivial zeros. See \Cref{v4-arith:rh:rem:trivial-zero}.
\end{remark}

\begin{proposition}[Bost--Connes ergodic decomposition on
\texorpdfstring{\(\mathcal{H}^{(p)}_{\mathrm{prim}}\)}{H\^{}(p)\_prim}]
\label{v4-arith:prim-zeta:prop:BC-ergodic}
\ClaimStatusProvedHere. The Bost--Connes
\(\mathbb{N}^{\times}\)-action on \(\mathcal{H}^{(p)}_{\mathrm{prim}}\)
via \(\mu_{n}\cdot e_{m,p}=e_{m+n,p}\) is ergodic, with KMS\(_{\beta}\)
states at inverse temperature \(\beta=s\) given by
\(\varphi_{s}(e_{n,p})=p^{-ns}\) and partition function
\(Z^{(p)}_{\mathrm{BC}}(s)=(1-p^{-s})^{-1}\) per prime.
\end{proposition}

\begin{proof}
The Bost--Connes Hecke partial isometries \(\mu_{n}\) act on the
single-oscillator basis \(\{e_{m,p}\}_{m\geq 1}\) by raising: \(\mu_{n}\)
sends \(e_{m,p}\) to \(e_{m+n,p}\) if compatible parity and to \(0\)
otherwise. The action is irreducible on
\(\mathcal{H}^{(p)}_{\mathrm{prim}}\) (single orbit generated by
\(e_{1,p}\) under \(\{\mu_{n}\}_{n\geq 1}\)), hence ergodic.

KMS condition: the time evolution \(\sigma_t\) of the Bost--Connes system
acts on Hecke operators by \(\sigma_t(\mu_n)=n^{it}\mu_n\). A state
\(\varphi\) on the Hecke algebra \(\mathbb{A}^{\mathrm{BC}}\) is a
KMS\(_\beta\) state if \(\varphi(AB)=\varphi(B\sigma_{i\beta}(A))\) for
analytic elements. The explicit KMS states at inverse temperature
\(\beta>1\) are \(\varphi_\beta(\mu_n\mu_m^*)=\delta_{n,m}n^{-\beta}/\zeta(\beta)\)
on the full algebra; see Bost--Connes 1995 Selecta Math.~1 Prop.~27 for
the complete verification. On the primary subspace
\(\mathcal{H}^{(p)}_\mathrm{prim}\), the restriction of this KMS state at
\(\beta=s\) assigns weight \(\varphi_s(e_{n,p})=p^{-ns}\) to each basis
vector \(e_{n,p}\) (this is the diagonal component of the Bost--Connes
KMS state restricted to the per-prime factor at inverse temperature
\(\beta=s\), with \(n^{-\beta}\) specialised to \(p^{-ns}\) for the prime-\(p\)
sector: Bost--Connes 1995 Prop.~27 specialised, not re-derived).

Partition function: the per-prime primary partition function is
\(Z^{(p)}_{\mathrm{BC}}(s)=\sum_{n\geq 1}\varphi_{s}(e_{n,p})=\sum_{n\geq 1}p^{-ns}=p^{-s}/(1-p^{-s})\).
The Euler factor \((1-p^{-s})^{-1}\) equals \(1+p^{-s}/(1-p^{-s})\), where the
\(1\) is the vacuum contribution \(\varphi_s(|0\rangle_p)=1\). The vacuum
\(|0\rangle_p\) is excluded from \(\mathcal{H}^{(p)}_\mathrm{prim}\) (Definition~\ref{v4-arith:prim-zeta:def:p-primary}) but contributes to
the KMS normalisation; in the restricted tensor product for the adelic
character, the per-prime Euler factor is counted as \((1-p^{-s})^{-1}\)
by the standard convention in which the KMS partition function includes
the vacuum unit, consistent with Bost--Connes 1995 Thm.~27's formula
\(Z_\mathrm{BC}(\beta)=\zeta(\beta)\) (which counts \(n=1,2,3,\ldots\)
including the unit \(\mu_1=\mathrm{id}\), contributing \(1^{-\beta}=1\)).
The stated partition function \(Z^{(p)}_\mathrm{BC}(s)=(1-p^{-s})^{-1}\)
thus uses the KMS convention, not the reduced character of
\(\mathcal{H}^{(p)}_\mathrm{prim}\) alone.
\end{proof}

\subsection{Archimedean Primary Subspace}
\label{v4-arith:prim-zeta:archimedean-primary}

\begin{definition}[archimedean primary Heisenberg]
\label{v4-arith:prim-zeta:def:arch-primary}
\ClaimStatusDefinitional. The archimedean Heisenberg factor
\(\mathcal{H}^{(\infty)}\) is the Fock representation of the rank-\(1\)
Heisenberg vertex algebra over \(\mathbb{R}\), with creation operators
\(\{\alpha_{-n}\}_{n\geq 1}\) and vacuum \(|0\rangle_{\infty}\). The
\emph{archimedean primary subspace}
\(\mathcal{H}^{(\infty)}_{\mathrm{prim}}\subset\mathcal{H}^{(\infty)}\)
is the subspace spanned by the Gaussian primary vector
\(|\text{Gauss}\rangle=e^{-\pi x^{2}}\cdot|0\rangle_{\infty}\) under the
coordinate realization of \(\mathcal{H}^{(\infty)}\) as
\(L^{2}(\mathbb{R})\) (Segal 1981 Prop.~3.2; Frenkel--Ben-Zvi 2004
\S3.5): concretely,
\begin{equation}
\label{v4-arith:prim-zeta:eq:arch-primary}
\mathcal{H}^{(\infty)}_{\mathrm{prim}}\;:=\;\mathbb{C}\cdot|\text{Gauss}\rangle\;=\;\mathbb{C}\cdot e^{-\pi x^{2}}|0\rangle_{\infty}.
\end{equation}
Descendant states \((\alpha_{-n})^{k}|\text{Gauss}\rangle\) for
\(n\geq 1,k\geq 1\) are excluded from the archimedean primary
subspace; these correspond to the Hermite-polynomial tower above the
Gaussian and carry the trivial-zero sector of \(\xi\).
\end{definition}

\begin{proposition}[archimedean primary character]
\label{v4-arith:prim-zeta:prop:arch-char}
\ClaimStatusProvedHere. Under the HY-zeta grading via conformal
weight times \(\log p_{\infty}\) (with archimedean ``prime'' frequency
\(\log p_{\infty}:=\text{archimedean analogue of } 1\) in the Tate
presentation; equivalently, \(L_{0}^{(\infty)}\) has continuous
spectrum \([0,\infty)\)), the character of
\(\mathcal{H}^{(\infty)}_{\mathrm{prim}}\) is
\begin{equation}
\label{v4-arith:prim-zeta:eq:arch-char}
\mathrm{char}\bigl(\mathcal{H}^{(\infty)}_{\mathrm{prim}}\bigr)(s)
\;=\;\int_{0}^{\infty}e^{-\pi x^{2}}\cdot x^{s-1}dx
\;=\;\tfrac{1}{2}\Gamma_{\mathbb{R}}(s)
\;=\;\tfrac{1}{2}\pi^{-s/2}\Gamma(s/2).
\end{equation}
The factor of \(2\) accommodates the archimedean sign involution
\(x\mapsto -x\): extending the integral over \(\mathbb{R}\) (not just
\((0,\infty)\)) gives \(\Gamma_{\mathbb{R}}(s)\) directly, which is the
Tate normalisation.
\end{proposition}

\begin{proof}
The Gaussian primary state \(|\text{Gauss}\rangle=e^{-\pi x^{2}}\)
under Mellin pairing against \(x^{s-1}\) on \((0,\infty)\):
\begin{equation}
\int_{0}^{\infty}e^{-\pi x^{2}}\cdot x^{s-1}dx
\;=\;\tfrac{1}{2}\pi^{-s/2}\Gamma(s/2),
\end{equation}
by substitution \(u=\pi x^{2}\), \(du=2\pi x\,dx\):
\begin{equation}
\tfrac{1}{2}\int_{0}^{\infty}e^{-u}\cdot(u/\pi)^{(s-1)/2}\cdot(u/\pi)^{-1/2}du/\pi^{1/2}
\;=\;\tfrac{1}{2}\pi^{-s/2}\int_{0}^{\infty}e^{-u}u^{s/2-1}du
\;=\;\tfrac{1}{2}\pi^{-s/2}\Gamma(s/2).
\end{equation}
Extension to \(\mathbb{R}\) doubles this: the full
Mellin--theta presentation is \(\int_{-\infty}^{\infty}e^{-\pi x^{2}}|x|^{s-1}dx=\pi^{-s/2}\Gamma(s/2)=\Gamma_{\mathbb{R}}(s)\)
(Tate 1950 \S2.5; Cassels--Fr\"ohlich 1967 \S13.3).
\end{proof}

\subsection{Adelic Primary Subspace}
\label{v4-arith:prim-zeta:adelic-primary}

\begin{definition}[primary subspace
\(\mathcal{V}^{\mathrm{prim}}_{\mathrm{Heis},\mathrm{prim}}\)]
\label{v4-arith:prim-zeta:def:primary-subspace}
\ClaimStatusDefinitional. The \emph{primary subspace} of the adelic
arithmetic Heisenberg Fock is the restricted tensor product
\begin{equation}
\label{v4-arith:prim-zeta:eq:adelic-primary}
\mathcal{V}^{\mathrm{prim}}_{\mathrm{Heis},\mathrm{prim}}\;:=\;
\mathcal{H}^{(\infty)}_{\mathrm{prim}}\otimes_{\mathbb{C}}\bigotimes_{p\text{ prime}}^{\prime}\mathcal{H}^{(p)}_{\mathrm{prim}}
\end{equation}
where:
\begin{itemize}
\item At \(\infty\), \(\mathcal{H}^{(\infty)}_{\mathrm{prim}}\) is the
Gaussian primary line
(\Cref{v4-arith:prim-zeta:def:arch-primary}).
\item At each finite \(p\), \(\mathcal{H}^{(p)}_{\mathrm{prim}}\) is
the single-oscillator primary space
(\Cref{v4-arith:prim-zeta:def:p-primary}).
\item The restricted tensor product is taken with respect to the
spherical vector \(e_{0,p}:=|0\rangle_{p}\) at almost all primes in
the limit inside the primary
(Tate 1950 \S2.3; Weil 1967 Ch.~VII \S1). More precisely, the
primary subspace extends the definition to include the trivial
``zero oscillator'' at almost all primes, but the restricted tensor
product converges in the \(L^{2}\)-adelic sense only after the
Bost--Connes normalisation factor per prime is absorbed into the
KMS state definition.
\end{itemize}
The underlying \(\Lambda\)-module structure at each finite \(p\) is
free of rank one
(the Iwasawa \(\mu\)-invariant is \(0\) on the single-oscillator
component, \Cref{v4-arith:pos:cor:mu-zero-heisenberg}).
\end{definition}

\begin{remark}[consistency with Virasoro bootstrap chapter]
\label{v4-arith:prim-zeta:rem:consistency-vb}
\Cref{v4-arith:vb:def:T_n} uses
\(\mathcal{V}^{\mathrm{prim}}_{\mathrm{Heis},\mathrm{prim}}\) as the
domain of the descendant operator \(T_{n}=(L_{0}^{\mathrm{Ar}})^{n}-(L_{0}^{\mathrm{Ar}}-\mathrm{id})^{n}\).
The present definition makes this operator well-defined: on
\(\mathcal{H}^{(p)}_{\mathrm{prim}}\), \(L_{0}^{\mathrm{Ar}}\) acts with
eigenvalue \(n\log p\) on \(e_{n,p}\); on
\(\mathcal{H}^{(\infty)}_{\mathrm{prim}}=\mathbb{C}\cdot|\text{Gauss}\rangle\),
\(L_{0}^{\mathrm{Ar}}\) acts as the Tate-archimedean \(s/2\) scaling
(distributional spectrum over \(\mathbb{R}\), regularised through
Mellin). The common domain for \(L_{0}^{\mathrm{Ar}}\) and its powers
is the algebraic tensor product of \(\bigoplus_{n\geq 1}\mathbb{C}e_{n,p}\)
with the Schwartz-archimedean Gaussian sector; spectral-calculus
extension is given in
\Cref{v4-arith:vb:def:T_n} under H-P\(^{\mathrm{Ar}}\).
\end{remark}

\begin{remark}[Virasoro reconciliation with
\Cref{v4-arith:rh:conj:primary-Virasoro}]
\label{v4-arith:prim-zeta:rem:Virasoro-reconciliation}
On the primary subspace
\(\mathcal{V}^{\mathrm{prim}}_{\mathrm{Heis},\mathrm{prim}}\), the
per-prime Virasoro action \(L_{m}^{(p)}\) is per-state finite: each
basis element \(e_{n,p}\) lies in a single per-prime Fock, and
\(L_{m}^{(p)}e_{n,p}=0\) for all \(m\neq 0\) (primary \(\Leftrightarrow\)
kernel of raising operators \(L_{m}^{(p)}\) for \(m\geq 1\), and
lowering \(L_{m}^{(p)}\) for \(m\leq -1\) produce descendants which are
excluded from the primary subspace). Hence
\(L_{m}^{\mathrm{Ar}}=\sum_{p}L_{m}^{(p)}\) acts coherently, and the
central-term sum \(\sum_{p}c^{(p)}\) regularises to \(c_{\mathrm{reg}}=2\)
(\Cref{v4-arith:rh:conj:primary-Virasoro}); the value \(c^{\mathrm{Ar}}=2\)
is established in \Cref{v4-arith:thm:c-as-residue} as the residue of
the Arakelov zeta function at \(s=1\), not via regularisation of
\(\sum_p p^{-s}\) at \(s=0\)
(the prime zeta function \(P(s)=\sum_{p}p^{-s}\) has a branch point at
\(s=0\) and no finite continuation, so direct regularisation of
\(\sum_p p^{0}\) is not available; the residue-at-\(s=1\) route through
\Cref{v4-arith:thm:c-as-residue} is the proof of record).
\end{remark}

\section{Character on Primary via Mellin--Theta}
\label{v4-arith:prim-zeta:character-primary}

\subsection{Finite-Place Primary Character}
\label{v4-arith:prim-zeta:finite-char}

\begin{theorem}[finite-place primary character]
\label{v4-arith:prim-zeta:thm:finite-char}
\ClaimStatusProvedHere. The character of
\(\bigotimes'_{p}\mathcal{H}^{(p)}_{\mathrm{prim}}\) under the
HY-zeta grading is the Dirichlet series \(\zeta(s)\), convergent
absolutely for \(\mathrm{Re}(s)>1\) and admitting meromorphic
continuation to \(\mathbb{C}\) with a simple pole at \(s=1\) and no
other poles:
\begin{equation}
\label{v4-arith:prim-zeta:eq:finite-char}
\mathrm{char}\Bigl(\bigotimes_{p\text{ prime}}^{\prime}\mathcal{H}^{(p)}_{\mathrm{prim}}\Bigr)(s)
\;=\;\prod_{p\text{ prime}}\bigl(1-p^{-s}\bigr)^{-1}
\;=\;\zeta(s),
\qquad\mathrm{Re}(s)>1.
\end{equation}
\end{theorem}

\begin{proof}
Per prime: the Bost--Connes KMS partition function on
\(\mathcal{H}^{(p)}_{\mathrm{prim}}\) at inverse temperature \(\beta=s\)
is \(Z^{(p)}_\mathrm{BC}(s)=(1-p^{-s})^{-1}\)
(\Cref{v4-arith:prim-zeta:prop:BC-ergodic}), where the Euler factor
\((1-p^{-s})^{-1}\) counts the KMS vacuum unit plus the single-oscillator
basis \(\{e_{n,p}\}_{n\geq 1}\) with weights \(p^{-ns}\). The restricted
tensor product of per-prime KMS partition functions produces
\(\prod_{p}(1-p^{-s})^{-1}\), the Euler product for
\(\zeta(s)\) (Riemann 1859; Cassels--Fr\"ohlich 1967 \S13.2).

Convergence: the product \(\prod_p(1-p^{-s})^{-1}\) converges absolutely for
\(\mathrm{Re}(s)>1\) by the bound \(\log(1-p^{-s})^{-1}=\sum_{k\geq 1}p^{-ks}/k\leq 2p^{-s}\)
for \(\mathrm{Re}(s)\geq 1\), and \(\sum_p p^{-\sigma}\) converges for
\(\sigma=\mathrm{Re}(s)>1\) (this is the prime zeta function \(P(\sigma)\),
finite for \(\sigma>1\); Landau 1909). Meromorphic continuation: the identity
\(\zeta(s)=(1-2^{1-s})^{-1}\sum_{n\geq 1}(-1)^{n-1}n^{-s}\)
(alternating zeta) gives continuation to \(\mathrm{Re}(s)>0\), then
functional equation gives continuation to \(\mathbb{C}\). The simple
pole at \(s=1\) is a standard computation
(\Cref{v4-arith:thm:c-as-residue}: \(\mathrm{Res}_{s=1}\zeta(s)=1\)).

Julia's 1990 primon gas identification
(\Cref{v4-arith:thm:P1-resolution} Step 1): the per-prime
single-oscillator Fock partition function equals
\((1-p^{-\beta})^{-1}\) at inverse temperature \(\beta\), and the
primon gas total partition is \(\zeta(\beta)\), identifying the
HY-zeta character with the Dirichlet series identically.
\end{proof}

\subsection{Archimedean Primary Character}
\label{v4-arith:prim-zeta:arch-char}

\begin{theorem}[archimedean primary character]
\label{v4-arith:prim-zeta:thm:arch-char}
\ClaimStatusProvedHere. The character of the archimedean primary
subspace \(\mathcal{H}^{(\infty)}_{\mathrm{prim}}\) is
\begin{equation}
\label{v4-arith:prim-zeta:eq:full-arch-char}
\mathrm{char}\bigl(\mathcal{H}^{(\infty)}_{\mathrm{prim}}\bigr)(s)
\;=\;\int_{\mathbb{R}}e^{-\pi x^{2}}|x|^{s-1}dx
\;=\;\pi^{-s/2}\Gamma(s/2)
\;=\;\Gamma_{\mathbb{R}}(s),
\end{equation}
the archimedean local \(L\)-factor of the trivial Dirichlet character.
\end{theorem}

\begin{proof}
\Cref{v4-arith:prim-zeta:prop:arch-char}: the Mellin pairing of the
Gaussian primary state with \(|x|^{s-1}\) gives \(\Gamma_{\mathbb{R}}(s)\).
The Mellin--theta machinery of Tate 1950 \S4.4 identifies this with
the archimedean local \(L\)-factor: \(L_{\infty}(s,\chi_{\mathrm{triv}})=\Gamma_{\mathbb{R}}(s)\),
with the Gaussian vector being the unique (up to scale) Tate test
function at \(\infty\) whose Mellin transform gives \(\Gamma_{\mathbb{R}}\).
\end{proof}

\subsection{Adelic Character: \(\xi(s)\)}
\label{v4-arith:prim-zeta:adelic-char}

\begin{theorem}[primitive zeta-character identity]
\label{v4-arith:prim-zeta:thm:adelic-char}
\ClaimStatusProvedHere. The character of the primary subspace
\(\mathcal{V}^{\mathrm{prim}}_{\mathrm{Heis},\mathrm{prim}}\) under the
adelic HY-zeta grading is the completed Riemann zeta function:
\begin{equation}
\label{v4-arith:prim-zeta:eq:adelic-char}
\mathrm{char}\bigl(\mathcal{V}^{\mathrm{prim}}_{\mathrm{Heis},\mathrm{prim}}\bigr)(s)
\;=\;\Gamma_{\mathbb{R}}(s)\cdot\zeta(s)\;=\;\xi(s).
\end{equation}
\end{theorem}

\begin{proof}
Factorise the adelic primary as
\(\mathcal{V}^{\mathrm{prim}}_{\mathrm{Heis},\mathrm{prim}}=\mathcal{H}^{(\infty)}_{\mathrm{prim}}\otimes_{\mathbb{C}}\bigotimes'_{p}\mathcal{H}^{(p)}_{\mathrm{prim}}\)
(\Cref{v4-arith:prim-zeta:def:primary-subspace}). The character of
a tensor product is the product of characters:
\begin{equation}
\mathrm{char}(\mathcal{V}^{\mathrm{prim}}_{\mathrm{Heis},\mathrm{prim}})(s)
\;=\;\mathrm{char}(\mathcal{H}^{(\infty)}_{\mathrm{prim}})(s)\cdot\mathrm{char}\bigl(\bigotimes'_{p}\mathcal{H}^{(p)}_{\mathrm{prim}}\bigr)(s).
\end{equation}
By \Cref{v4-arith:prim-zeta:thm:arch-char},
\(\mathrm{char}(\mathcal{H}^{(\infty)}_{\mathrm{prim}})(s)=\Gamma_{\mathbb{R}}(s)\).
By \Cref{v4-arith:prim-zeta:thm:finite-char},
\(\mathrm{char}(\bigotimes'_{p}\mathcal{H}^{(p)}_{\mathrm{prim}})(s)=\zeta(s)\).
Product: \(\Gamma_{\mathbb{R}}(s)\zeta(s)=\xi(s)\), the standard
completion by Tate 1950 \S4.4 / Cassels--Fr\"ohlich 1967 \S13.4.
\end{proof}

\begin{corollary}[primary character of Heisenberg is \(\xi\)]
\label{v4-arith:prim-zeta:cor:primary-is-xi}
\ClaimStatusProvedHere.
\(\mathrm{char}(\mathcal{V}^{\mathrm{prim}}_{\mathrm{Heis},\mathrm{prim}})(s)=\xi(s)\),
a meromorphic function of \(s\) with simple poles only at \(s=0\) and
\(s=1\) (both residues \(\pm 1\)), and with functional equation
\(\xi(s)=\xi(1-s)\) (Riemann 1859; Tate 1950).
\end{corollary}

\begin{proof}
Combine \Cref{v4-arith:prim-zeta:thm:adelic-char} with Tate's
functional equation
(\Cref{v4-arith:thm:P2-resolution}(3)).
\end{proof}

\section{Koszul-Dual Character (Bost--Connes Side)}
\label{v4-arith:prim-zeta:koszul-dual-char}

\subsection{Cobar of Primary Heisenberg}
\label{v4-arith:prim-zeta:cobar-primary}

\begin{definition}[cobar algebra on primary]
\label{v4-arith:prim-zeta:def:cobar-primary}
\ClaimStatusDefinitional. Let \(C_{\mathrm{prim}}:=B^{\mathrm{ch,Ar}}(\mathcal{V}^{\mathrm{prim}}_{\mathrm{Heis},\mathrm{prim}})\)
be the arithmetic bar coalgebra on the primary subspace, using the
bar construction of \Cref{v4-arith:def:arith-bar-concrete}
restricted to primary generators. The \emph{cobar algebra on primary}
is \(A_{\mathrm{prim}}:=(C_{\mathrm{prim}})^{\vee}=\Omega^{\mathrm{ch,Ar}}(\mathcal{V}^{\mathrm{prim}}_{\mathrm{Heis},\mathrm{prim}})\).
By \Cref{v4-arith:cor:koszul-dual-BC}, \(A_{\mathrm{prim}}\) is the
primary-sector projection of the Bost--Connes Hecke algebra
\(\mathbb{A}^{\mathrm{BC}}\), generated by the Hecke operators
\(\{\mu_{n}\}_{n\geq 1}\) paired with the primary oscillators
\(\{e_{n,p}\}\) via \(\mu_{p^{n}}\leftrightarrow e_{n,p}\)
(\Cref{v4-arith:pos:thm:weight-comp}(3)).
\end{definition}

\subsection{Bost--Connes KMS Partition on Primary}
\label{v4-arith:prim-zeta:BC-KMS}

\begin{theorem}[Bost--Connes partition on primary]
\label{v4-arith:prim-zeta:thm:BC-partition}
\ClaimStatusProvedHere. The KMS\(_{\beta}\) partition function of the
cobar algebra \(A_{\mathrm{prim}}\) at inverse temperature
\(\beta=s\) is
\begin{equation}
\label{v4-arith:prim-zeta:eq:BC-partition}
Z_{A_{\mathrm{prim}}}(\beta)\;=\;\zeta(\beta),\qquad \mathrm{Re}(\beta)>1.
\end{equation}
For \(\mathrm{Re}(\beta)\leq 1\), \(Z_{A_{\mathrm{prim}}}(\beta)\)
admits meromorphic continuation via Bost--Connes 1995 Selecta Math.~1
Thm.~27 plus analytic continuation of \(\zeta\).
\end{theorem}

\begin{proof}
Bost--Connes 1995 Thm.~27: the full Bost--Connes Hecke algebra
\(\mathbb{A}^{\mathrm{BC}}\) has KMS\(_{\beta}\) partition function
\(\zeta(\beta)\) at inverse temperature \(\beta>1\). Restriction to the
primary sector: the primary cobar \(A_{\mathrm{prim}}\) is generated
by \(\{\mu_{n}\}\) with single-Adams weight \(-\log n\) (no higher-Adams
contribution beyond the Bost--Connes single-isometry generators);
its KMS\(_{\beta}\) partition function is the full Bost--Connes
partition restricted to \(A_{\mathrm{prim}}\), which is
\(\zeta(\beta)\) identically because the primary sector carries the
full Dirichlet-series signal. The higher-Fock (multi-oscillator)
sectors contribute additional \(\zeta(n\beta)\) factors for
\(n\geq 2\) that are \emph{not} in the primary.
\end{proof}

\subsection{Archimedean Factor via Fourier Duality}
\label{v4-arith:prim-zeta:arch-fourier}

\begin{proposition}[archimedean character on cobar]
\label{v4-arith:prim-zeta:prop:arch-cobar}
\ClaimStatusProvedHere. The archimedean factor of \(A_{\mathrm{prim}}^{!}\)
has character
\begin{equation}
\label{v4-arith:prim-zeta:eq:arch-cobar}
\mathrm{char}\bigl((\mathcal{H}^{(\infty)}_{\mathrm{prim}})^{!}\bigr)(s)
\;=\;\Gamma_{\mathbb{R}}(1-s)
\;=\;\pi^{-(1-s)/2}\Gamma((1-s)/2),
\end{equation}
the Fourier dual of the archimedean primary character.
\end{proposition}

\begin{proof}
Koszul duality on the archimedean sector is classical Positselski
over \(\mathbb{R}\) (\Cref{v4-arith:pos:thm:gaussian-positselski}):
the bar--cobar pairing identifies the Gaussian primary vector with
its Fourier dual
\(\mathcal{F}_{\mathbb{R}}e^{-\pi x^{2}}=e^{-\pi y^{2}}\). The
character of \((\mathcal{H}^{(\infty)}_{\mathrm{prim}})^{!}\) is the
Mellin pairing of \(e^{-\pi y^{2}}\) against \(|y|^{(1-s)-1}\) on
\(\mathbb{R}\), i.e.\
\begin{equation}
\int_{\mathbb{R}}e^{-\pi y^{2}}|y|^{(1-s)-1}dy\;=\;\pi^{-(1-s)/2}\Gamma((1-s)/2)\;=\;\Gamma_{\mathbb{R}}(1-s).
\end{equation}
The substitution \(s\mapsto 1-s\) is the Koszul-dual / Tate-Fourier
transform on the archimedean component; equivalently, Mellin pairing
against a shifted exponent.
\end{proof}

\subsection{Adelic Koszul-Dual Character}
\label{v4-arith:prim-zeta:koszul-dual-adelic}

\begin{theorem}[Koszul-dual adelic character]
\label{v4-arith:prim-zeta:thm:koszul-dual-adelic}
\ClaimStatusProvedHere. The character of the Koszul-dual
\((\mathcal{V}^{\mathrm{prim}}_{\mathrm{Heis},\mathrm{prim}})^{!}\)
under the adelic HY-zeta grading, with the substitution
\(s\mapsto 1-s\) at the archimedean place, is
\begin{equation}
\label{v4-arith:prim-zeta:eq:koszul-dual-adelic}
\mathrm{char}\bigl((\mathcal{V}^{\mathrm{prim}}_{\mathrm{Heis},\mathrm{prim}})^{!}\bigr)(s)
\;=\;\Gamma_{\mathbb{R}}(1-s)\cdot\zeta(1-s)\;=\;\xi(1-s).
\end{equation}
\end{theorem}

\begin{proof}
Factorise the Koszul-dual adelic primary. The finite-place Koszul
dual is the Bost--Connes cobar with partition function
\(\zeta(s)\) (\Cref{v4-arith:prim-zeta:thm:BC-partition}), but under
the Koszul-dual substitution \(\beta\mapsto 1-\beta\) this becomes
\(\zeta(1-s)\): the Koszul-dual character is the Bost--Connes
partition at inverse temperature \(\beta=1-s\), not \(\beta=s\). The
archimedean factor is \(\Gamma_{\mathbb{R}}(1-s)\) by
\Cref{v4-arith:prim-zeta:prop:arch-cobar}. Product:
\(\Gamma_{\mathbb{R}}(1-s)\cdot\zeta(1-s)=\xi(1-s)\).

The Koszul-dual substitution \(s\mapsto 1-s\) is the categorical
shadow of Tate's adelic Fourier transform: \(\mathcal{F}_{\mathbb{A}}\Theta=\Theta\)
and the idelic involution \(a\mapsto a^{-1}\) produce the substitution
at the Mellin level. See \Cref{v4-arith:pos:rem:adelic-fe} for the
bar--cobar shadow of Fourier self-duality.
\end{proof}

\section{Primitive Zeta-Character Theorem}
\label{v4-arith:prim-zeta:main-theorem}

\subsection{Assembly}
\label{v4-arith:prim-zeta:assembly}

\begin{proof}[Proof of \Cref{v4-arith:prim-zeta:thm:main}]
The character identity
(\ref{v4-arith:prim-zeta:eq:main})
\begin{equation*}
\mathrm{char}(\mathcal{V}^{\mathrm{prim}}_{\mathrm{Heis},\mathrm{prim}})(s)
\;=\;\Gamma_{\mathbb{R}}(s)\,\zeta(s)\;=\;\xi(s)
\end{equation*}
is \Cref{v4-arith:prim-zeta:thm:adelic-char}: factorise the adelic
primary into finite-place and archimedean components, apply
\Cref{v4-arith:prim-zeta:thm:finite-char} (Julia 1990 primon gas
+ Euler product) and
\Cref{v4-arith:prim-zeta:thm:arch-char} (Tate 1950 Mellin--theta on
Gaussian), tensor to obtain
\(\Gamma_{\mathbb{R}}(s)\zeta(s)=\xi(s)\).

The character identity
(\ref{v4-arith:prim-zeta:eq:koszul-dual})
is \Cref{v4-arith:prim-zeta:thm:koszul-dual-adelic}: the Koszul-dual
character satisfies \(\mathrm{char}_{A^{!}}(s)=\xi(1-s)\), so
\(\mathrm{char}_{C}(s)=\mathrm{char}_{A^{!}}(1-s)=\xi(1-s)\big|_{1-s\mapsto s}=\xi(s)\).
Tate's functional equation \(\xi(s)=\xi(1-s)\)
(\Cref{v4-arith:thm:P2-resolution}(3)) then gives
\(\mathrm{char}_{C}(s)=\xi(s)=\xi(1-s)=\mathrm{char}_{A^{!}}(s)\)
as an identity of meromorphic functions.

The character identity holds unconditionally; the \emph{pairing}
(not just the character identity) between \(C\) and \(A^{!}\) requires
the four remaining open Positselski sub-items of
\Cref{v4-arith:pos:rem:open-sub-items} (continuous duals; derived
completion; faithful lifting; Koszul-dual Bost--Connes), which
this chapter does not supply. The primitive zeta-character theorem
\emph{is} the character-level identity; the categorical identity
remains conditional.
\end{proof}

\subsection{Disjointness of Sources}
\label{v4-arith:prim-zeta:disjoint-sources}

\begin{remark}[sources used in the proof]
\label{v4-arith:prim-zeta:rem:sources}
The four verifications use four disjoint primary sources:
\begin{enumerate}
\item \textbf{Tate 1950 Mellin--theta}
(\Cref{v4-arith:cl:P2-resolution},
Tate \emph{Fourier Analysis in Number Fields and Hecke's
Zeta-Functions}, Princeton 1950 thesis reprinted Cassels--Fr\"ohlich
1967 \S13): archimedean \(\Gamma_{\mathbb{R}}(s)\) computation.
\item \textbf{Julia 1990 primon gas}
(B.~Julia, \emph{Statistical theory of numbers}, in Number Theory
and Physics, Springer Proc.\ 47, 1990): finite-place
\((1-p^{-s})^{-1}\) per prime and Euler product to \(\zeta(s)\).
\item \textbf{Bost--Connes 1995}
(J.--B.~Bost and A.~Connes, \emph{Hecke algebras, type III factors
and phase transitions with spontaneous symmetry breaking in number
theory}, Selecta Math.~(N.S.)~1 (1995):411--457): KMS\(_{\beta}\)
partition function \(\zeta(\beta)\) on the Bost--Connes Hecke cobar,
and Koszul-dual substitution \(\beta\mapsto 1-\beta\) reading.
\item \textbf{Positselski 2011 classical character pairing}
(L.~Positselski, \emph{Two kinds of derived categories, Koszul
duality, and comodule--contramodule correspondence}, Mem.~AMS 996
(2011), arXiv:0905.2621, Thm.~5.2): classical Koszul pairing at
character level over \(\mathbb{C}\), used to identify the tensor
factors on the archimedean side via Gaussian Positselski
(\Cref{v4-arith:pos:thm:gaussian-positselski}).
\end{enumerate}
All four predate the arithmetic bar--cobar programme, with no internal
cross-dependency in that literature. Each input
is drawn from a different mathematical tradition (analytic number
theory; statistical physics; non-commutative geometry; categorical
algebra), which gives the disjointness of
\Cref{v4-arith:prim-zeta:rem:sources-disjoint} its substance.
\end{remark}

\begin{lemma}[source disjointness audit]
\label{v4-arith:prim-zeta:rem:sources-disjoint}
\ClaimStatusProvedHere. The four sources Tate 1950, Julia 1990,
Bost--Connes 1995, and Positselski 2011 have no cross-citation
relation among them. Specifically:
\begin{itemize}
\item Tate 1950 (analytic number theory thesis) pre-dates all
others and cites no bar--cobar or operator-algebraic sources.
\item Julia 1990 (statistical physics proceedings) cites Riemann
and classical analytic number theory; does not cite Tate 1950
thesis framework, Positselski, or Bost--Connes.
\item Bost--Connes 1995 cites Tate 1950 for adelic machinery but
treats it as black-box input; does not cite Positselski 2011 (which
post-dates it) or Julia 1990 (which it parallels but does not cite).
\item Positselski 2011 cites neither Tate 1950, Julia 1990, nor
Bost--Connes 1995; works in pure categorical algebra.
\end{itemize}
No pair among the four sources shares a citation chain
through any intermediate. This satisfies the Vol IV disjoint-source
criterion for realization decorators.
\end{lemma}

\section{Consequence for P1}
\label{v4-arith:prim-zeta:consequence-p1}

\subsection{Reduction of Open Sub-Items}
\label{v4-arith:prim-zeta:reduction-sub-items}

\begin{corollary}[P1 conditional residual reduced from five to four]
\label{v4-arith:prim-zeta:cor:p1-residual-reduced}
\ClaimStatusProvedHere. With
\Cref{v4-arith:prim-zeta:thm:main} inscribed, the primitive
zeta-character sub-item of
\Cref{v4-arith:pos:rem:open-sub-items}(4) is closed. The remaining
open Positselski sub-items are:
\begin{enumerate}
\item[(1)] Completed continuous duals (\Cref{v4-arith:pos:rem:open-sub-items}(1)).
\item[(2)] Derived completion (\Cref{v4-arith:pos:rem:open-sub-items}(2)).
\item[(3)] Faithful lifting along Nakayama (\Cref{v4-arith:pos:rem:open-sub-items}(3)).
\item[(5)] Koszul-dual Bost--Connes presentation (\Cref{v4-arith:pos:rem:open-sub-items}(5)).
\end{enumerate}
The arith.\ Positselski input on the four-input conditional spine
(\Cref{v4-arith:w3-audit-B:thm:four-input-spine}) is conditionalised
on these four sub-items, not five.
\end{corollary}

\begin{proof}
\Cref{v4-arith:prim-zeta:thm:main} establishes the primitive
zeta-character theorem as a standalone theorem, closing sub-item
(4) of \Cref{v4-arith:pos:rem:open-sub-items}. The other four
sub-items are independent and remain open.
\end{proof}

\begin{remark}[status of P1 biconditional after this chapter]
\label{v4-arith:prim-zeta:rem:p1-status}
\Cref{v4-arith:pos:thm:P1-full-unconditional} carries conditional
status on the five Positselski sub-items. After this chapter:
\begin{itemize}
\item Sub-item (4) [primitive zeta-character] is closed.
\item Sub-items (1), (2), (3), (5) remain open.
\end{itemize}
The status of
\Cref{v4-arith:pos:thm:P1-full-unconditional} accordingly tightens
from ``conditional on five'' to ``conditional on four'' (continuous
duals, derived completion, faithful lifting, Koszul-dual BC). The
load-bearing theorem
\Cref{v4-arith:thm:G-row-Riemann-functional-equation} remains
\emph{Theorem-waiting}: forward direction unconditional
(\Cref{v4-arith:cl:cor:P1-unconditional-forward}), full biconditional
conditional on four sub-items (down from five).
\end{remark}

\subsection{Character-Level Fingerprint of the Functional Equation}
\label{v4-arith:prim-zeta:fingerprint}

\begin{proposition}[character-level Riemann functional equation]
\label{v4-arith:prim-zeta:prop:char-functional-eq}
\ClaimStatusProvedHere. At the level of characters of the primary
subspace and its Koszul dual,
\begin{equation}
\label{v4-arith:prim-zeta:eq:char-functional-eq}
\mathrm{char}\bigl(\mathcal{V}^{\mathrm{prim}}_{\mathrm{Heis},\mathrm{prim}}\bigr)(s)
\;=\;\mathrm{char}\bigl((\mathcal{V}^{\mathrm{prim}}_{\mathrm{Heis},\mathrm{prim}})^{!}\bigr)(s)
\;\Longleftrightarrow\;
\xi(s)\;=\;\xi(1-s).
\end{equation}
\end{proposition}

\begin{proof}
By \Cref{v4-arith:prim-zeta:thm:adelic-char} and
\Cref{v4-arith:prim-zeta:thm:koszul-dual-adelic}, the two characters
equal \(\xi(s)\) and \(\xi(1-s)\) respectively. Equality of characters
is therefore the functional equation. Tate 1950 establishes the
implication \(\xi(s)=\xi(1-s)\) unconditionally
(\Cref{v4-arith:thm:P2-resolution}(3)); the character equality is
the Fock-operator-theoretic fingerprint of the Riemann functional
equation, inscribed on the primary subspace of the arithmetic
Heisenberg.
\end{proof}

\begin{remark}[what this proposition does not say]
\label{v4-arith:prim-zeta:rem:caveat-fingerprint}
\Cref{v4-arith:prim-zeta:prop:char-functional-eq} asserts that the
character identity and Tate's functional equation are equivalent
at the level of characters of the primary subspace. It does
\emph{not} prove the Riemann Hypothesis, nor does it close the P1
biconditional at the level of bar--cobar pairings (this requires
the Positselski equivalence, conditional on the four remaining
open sub-items).

The proposition is the character-level analogue of
\Cref{v4-arith:thm:G-row-Riemann-functional-equation}: Riemann
functional equation as the character-level shadow of Koszul
self-duality on the primary subspace. Full biconditional at the
pairing level requires arithmetic Positselski, which is not
closed by this chapter.
\end{remark}

\section{Remaining Four Positselski Sub-Items}
\label{v4-arith:prim-zeta:remaining-four}

For reader navigation, the four remaining open Positselski sub-items
are restated with explicit pointers to where they are expected to be
addressed in subsequent arithmetic-branch chapters.

\begin{description}
\item[\emph{(1) Completed continuous duals.}] Pontryagin-dual continuous
duals of completed Iwasawa modules (Brumer 1966 \S1--2;
Coates--Sujatha 2006 \S1.7), compatible with the Nekov\'a\v{r}
Selmer-complex formalism of \Cref{v4-arith:pos:att-2-co-contra}.
Expected closure: continuous-duality chapter combining Brumer with
Nekov\'a\v{r} 2006 Ast\'erisque 310.
\item[\emph{(2) Derived completion.}] Derived \(\mathfrak{m}\)-adic
completion of comodule and contramodule resolutions commuting with
base change \(\Lambda\to\mathbb{F}_{p}[[T]]\). Positselski 2018
arXiv:1209.2995 \S4.6 covers the Artinian base case; the Weierstrass
inverse limit is the open sub-item. Expected closure: Coates--Sujatha
2006 inverse-limit plus Positselski 2018 Artinian plus
Mittag-Leffler.
\item[\emph{(3) Faithful lifting along Nakayama.}] Essential
surjectivity of the base-change functor
\(\mathrm{D}^{\mathrm{co}}_{\Lambda}\to\mathrm{D}^{\mathrm{co}}_{\mathbb{F}_{p}[[T]]}\)
under \(\mu=0\) (\Cref{v4-arith:pos:prop:lifting}). Requires sub-item
(2) plus independent obstruction-theory check. Expected closure:
Nakayama-lifting chapter combining Illusie 1971 with Bhatt--Scholze
derived completion.
\item[\emph{(5) Koszul-dual Bost--Connes presentation.}] The
identification of the cobar algebra \(A=C^{\vee}\) with the
Bost--Connes Hecke algebra \(\mathbb{A}^{\mathrm{BC}}\)
(\Cref{v4-arith:cor:koszul-dual-BC}; currently
\texttt{ClaimStatusConjectured}). Requires independent
cross-verification that the Bost--Connes presentation carries the
Positselski pairing structure
(\Cref{v4-arith:pos:thm:weight-comp}). Expected closure:
Bost--Connes--Koszul chapter combining Bost--Connes 1995 \S1.2 with
the twisting-cochain formalism of Lefevre-Hasegawa 2003
(\Cref{v4-arith:pos:thm:lambda-bar-cobar}).
\end{description}

\subsection{Residual Ledger}
\label{v4-arith:prim-zeta:residual-ledger}

\begin{table}[h]
\centering
\small
\begin{tabular}{c|p{6cm}|l}
Sub-item & Content & Status after this chapter \\
\hline
(1) & Completed continuous duals (Brumer, Nekov\'a\v{r}) & Open \\
(2) & Derived completion (Positselski 2018, Weierstrass tower) & Open \\
(3) & Faithful lifting along Nakayama (\(\mu=0\), obstruction) & Open \\
(4) & \textbf{Primitive zeta-character theorem} & \textbf{Closed (this chapter)} \\
(5) & Koszul-dual Bost--Connes presentation & Open \\
\end{tabular}
\caption{Status of the five open Positselski sub-items of
\Cref{v4-arith:pos:rem:open-sub-items} after \Cref{v4-arith:prim-zeta:thm:main}.
Sub-item (4) is closed by
\Cref{v4-arith:prim-zeta:thm:main}; four remain open.}
\label{v4-arith:prim-zeta:tab:residual}
\end{table}

\section{Summary}
\label{v4-arith:prim-zeta:summary}

\begin{enumerate}
\item The \emph{primary subspace}
\(\mathcal{V}^{\mathrm{prim}}_{\mathrm{Heis},\mathrm{prim}}\)
(\Cref{v4-arith:prim-zeta:def:primary-subspace}) is defined as the
adelic tensor product of the per-prime single-oscillator primary
\(\mathcal{H}^{(p)}_{\mathrm{prim}}\) and the archimedean Gaussian
primary \(\mathcal{H}^{(\infty)}_{\mathrm{prim}}\). The definition
is consistent with the usage in
Ch.~\ref{v4-arith:virasoro-bootstrap} (Virasoro positivity) and
Ch.~\ref{v4-arith:rh-reformulation} (Li criterion).
\item \Cref{v4-arith:prim-zeta:thm:main} (primitive zeta-character
theorem): under the HY-zeta grading,
\(\mathrm{char}(\mathcal{V}^{\mathrm{prim}}_{\mathrm{Heis},\mathrm{prim}})(s)=\xi(s)\),
unconditionally.
\item The four verifications
(\Cref{v4-arith:prim-zeta:ledger}) use Tate 1950 (archimedean),
Julia 1990 (finite places), Bost--Connes 1995 (Koszul-dual cobar
side), and Positselski 2011 (classical character pairing); the four
sources are pairwise disjoint
(\Cref{v4-arith:prim-zeta:rem:sources-disjoint}).
\item \Cref{v4-arith:prim-zeta:prop:char-functional-eq}: at the
character level of the primary subspace, Koszul self-duality
\(\mathrm{char}_{C}(s)=\mathrm{char}_{A^{!}}(s)\) is equivalent to the
Riemann functional equation \(\xi(s)=\xi(1-s)\). This is the primary
character-level fingerprint of the load-bearing theorem
(\Cref{v4-arith:thm:G-row-Riemann-functional-equation}).
\item \Cref{v4-arith:prim-zeta:cor:p1-residual-reduced}: sub-item (4) of
\Cref{v4-arith:pos:rem:open-sub-items} [primitive zeta-character] is closed; four Positselski
sub-items remain open. The arith.\ Positselski input on the
four-input conditional spine
(\Cref{v4-arith:w3-audit-B:thm:four-input-spine}) is conditionalised
on four open sub-items, not five.
\item The load-bearing theorem status is preserved:
\emph{Theorem-waiting}, with forward direction unconditional and
full biconditional conditional on four (not five) Positselski
sub-items.
\end{enumerate}

This chapter closes sub-item (4) of the five open Positselski
sub-items of \Cref{v4-arith:thm:arithmetic-positselski}. The
character-level primitive zeta-character identity is a theorem, proved
unconditionally via Tate's Mellin--theta. The categorical pairing at
the bar--cobar level remains conditional on the four remaining open
sub-items. The arithmetic Positselski biconditional is conditionalised
on four (not five) sub-items.
